\documentclass[11pt,letterpaper]{article}
\usepackage[margin=0.85in]{geometry}
\usepackage[T1]{fontenc}
\usepackage{lmodern,microtype}
\usepackage{amsmath,amssymb,amsthm,booktabs,longtable,array,xcolor}
\usepackage{enumitem,needspace}
\usepackage{hyperref}
\usepackage{fancyhdr}
\definecolor{inkblue}{RGB}{29,55,82}
\hypersetup{colorlinks=true,linkcolor=inkblue,urlcolor=inkblue,citecolor=inkblue,
 pdftitle={PA194QNS: construction, comparison, and preservation proof},
 pdfauthor={Research construction and audit}}
\pagestyle{fancy}\fancyhf{}
\fancyhead[L]{\small PA194QNS}\fancyhead[R]{\small Construction and preservation}
\fancyfoot[C]{\thepage}\setlength{\headheight}{14pt}
\setlength{\emergencystretch}{3em}
\setlist{itemsep=3pt,topsep=5pt}
\newtheorem{theorem}{Theorem}[section]
\newtheorem{lemma}[theorem]{Lemma}
\newtheorem{corollary}[theorem]{Corollary}
\newcommand{\PA}{\mathsf{PA}}
\newcommand{\Con}{\operatorname{Con}}
\newcommand{\Halt}{\operatorname{Halt}}
\newcommand{\Prf}{\operatorname{Prf}}
\newcommand{\AX}{\mathsf{AX}}
\newcommand{\MP}{\mathsf{MP}}
\newcommand{\Shift}{\mathsf{S}}
\newcommand{\file}[1]{\path{#1}}
\title{\textcolor{inkblue}{PA194QNS}\\[4pt]
 A distinct 194-state machine for $\Con(\PA)$\\[8pt]
 \large Construction, comparison, and a finite-trace preservation proof}
\author{Research construction and audit}
\date{13 November 2026, Pacific time\\[3pt]
 \small ProveIt snapshot \texttt{317b07268681b44ff352d173a162e2a8c241b2fe}}
\begin{document}
\maketitle
\begin{abstract}
No machine with fewer than 194 working states was established. This report
supplies the requested fallback: PA194QNS, a complete blank-start binary
single-tape Turing machine with 194 working states, and the companion PA194QS.
The primary combines the supplied exchanged-premise axiom emitter, the
no-push successful modus-ponens protocol, and a symmetric implementation of
the original sum-based Shift guard. Its component ideas are credited to the
supplied research. The new result is their concrete compatible composition,
an exact all-register source-preservation proof, and independent physical
and corpus checks.

The main new preservation theorem is that PA proves the blank-halting
equivalence of PA194QNS and the supplied PA194QN. Composing it with the
retained ordinary mathematical PA calibration yields
$\PA\vdash\neg\Halt(\mathrm{PA194QNS})\leftrightarrow\Con(\PA)$.
The source has 3,651 primitive sites; its construction checker verifies
4,751 symbolic outcomes, all 32 simultaneous fold cases, and the entire
58,004-step blank initialization. The separate literal-tape audit rederives
the register-list lemmas from transition rows. These are written mathematical
proofs and executable certificates, not a newly compiled end-to-end Lean or
Rocq proof. The repository's separate iPRA formalization goal remains distinct.
\end{abstract}
\newpage
\tableofcontents
\newpage

\section{Result and the literal machine}

Let $H_M$ mean that the machine $M$ has a finite halting computation on its
specified blank input. Fix the usual arithmetization of PA proofs and set
\[
 \Con(\PA)\equiv\neg\exists p\,\Prf_{\PA}(p,\ulcorner0=1\urcorner).
\]
The result is a statement about a proof of an equivalence inside PA; it
does not assert that PA proves its own consistency or the machine's
nonhalting separately.

\begin{theorem}[Primary construction]
Using the supplied syntactic PA/calculus calibration detailed in
Section~\ref{sec:logical}, the explicit table PA194QNS satisfies
\[
 \PA\vdash H_{\mathrm{PA194QNS}}
 \leftrightarrow\exists p\,\Prf_{\PA}(p,\ulcorner0=1\urcorner).
\]
Consequently PA proves
$\neg H_{\mathrm{PA194QNS}}\leftrightarrow\Con(\PA)$.
\end{theorem}

The tape is two-way infinite, with alphabet $\{0,1\}$ and blank symbol zero.
Initially every cell is zero, the head is at coordinate zero, and the
control is zero. Working controls are $0,\ldots,193$; the terminal halt label
is 194 and has no transition rows. Each step writes one bit and moves one
cell left or right. There are 388 transitions, two for every working
control. No program, data, or initialization machine is supplied outside
those 194 controls.

The authoritative table is
\file{machines/PA194QNS/machine.tm}. Its row format is
\begin{center}\ttfamily state\quad read\quad write\quad L-or-R\quad next\end{center}
Its SHA-256 digest is
\begin{quote}\small
\nolinkurl{7eb36aea69e5de46cc47fa7aab0191f191657ab500dafb808dc3c394640591b8}.
\end{quote}
The native presentation uses start 74 and halt 254. A labels array includes
the exact dense/native bijection, including the halt label. Appendix~\ref{app:table}
prints every primary transition. All 194 working controls are reachable in
the finite directed control graph. This graph fact is separate from which
configurations are reachable on the blank tape.

% Generated by generate_includes.py.
\begin{center}\small
\begin{tabular}{lrr}
\toprule Verified quantity & PA194QNS & PA194QS\\\midrule
Working controls / transition rows & 194 / 388 & 194 / 388\\
Primitive source sites & 3,651 & 3,670\\
Zero/positive source outcomes & 4,751 & 4,776\\
Construction symbolic steps & 1,512,246 & 1,518,350\\
Independent literal symbolic steps & 1,106,186 & 1,110,647\\
Maximum literal steps per contract & 397 & 397\\
Smallest duration constant & 3 & 3\\
Generic list lemmas rederived & 4 & 4\\
Simultaneous local fold cases & 32 & 32\\
Literal blank initialization steps & 58,004 & 58,004\\
\bottomrule\end{tabular}\end{center}
The two symbolic step totals count operations in their respective proof
engines; they are not measurements of a particular natural-valued TM run.
Both engines cover arbitrary register values, not a finite register bound.


\paragraph{Companion.}
PA194QS keeps the original successful-MP dummy push and otherwise has the
same exchanged-premise emitter and symmetric original-sum guard. In native
labels it differs from QNS only at $(164,1)$: QS goes to 98, whereas QNS
goes to 114, in both cases writing 1 and moving right. Its complete dense
table and labels are included. The same comparator-preservation proof
applies, with the corresponding 3,673-site parent source. Its PA calibration
uses the original push protocol and the same exchanged-premise proof
translation.

\section{Comparison of the supplied corpus}
\label{sec:comparison}

The snapshot is pinned to the commit on the title page. It includes all
51 filed PA research packets and the three then-current incoming archives:
\file{PA194NS_research.zip}, \file{pa194haz_research_2026-11-14.zip}, and
\file{pa194qn_research_2026-11-13.zip}. The ledger in
Appendix~\ref{app:ledger} records every one of these 54 deliveries.
Larger designs, diagnostic tables, rejected contractions, duplicate
deliveries, and compiler intermediates are distinguished from accepted
manuscript constructions.

The repository's existing census records 4,984 originally delivered
\file{.tm} files, reduced to 490 table/machine classes when duplicate
state numberings were removed: 444 classes with a start and 46 unresolved
table-only classes. It reports 34 filed 194-state table/machine classes
(26 with a declared start and eight table-only classes). Those are
historical census quantities, not a count of independent consistency
theorems. The fresh audit operates on the current, already deduplicated
tree and recursively opens the surviving archives, so its occurrence
count is intentionally different.

% Generated by generate_includes.py.
\begin{center}\small
\begin{tabular}{lr}
\toprule Fresh audit quantity & Count\\\midrule
Delivery packets in the ledger & 54\\
Tracked candidate files matched to Git blobs & 2,431\\
Recognized literal table occurrences & 1,404\\
Distinct literal table byte strings & 722\\
Distinct ZIP payloads recursively opened & 17\\
Literal 194-state byte forms & 94\\
Whole-graph classes from literal tables & 39\\
JSON files additionally examined & 6,735\\
Embedded complete-table occurrences recovered & 62\\
Combined prior 194-state whole-graph classes & 40\\
\bottomrule\end{tabular}\end{center}
The JSON supplement recovers one extra class, PA194N-AD, that is absent
as a standalone literal table. Each of the 40 classes has an unchanged
raw representative and a source locator in the archive. Nine native
representatives do not declare a start; their null start annotations are
preserved. The novelty test is intentionally independent of that uncertainty.

The recognition scope is complete numeric five-column binary L/R tables
in literal table files and in JSON strings or arrays of five-column rows.
It does not interpret arbitrary source programs, arbitrary JSON transition
dictionaries, or every possible machine serialization. Thus the claim is
exhaustive for the recognized artifact formats in the pinned tree, rather
than a claim about every design that could be implicit in arbitrary code.


For novelty, the comparison uses full action tuples, rather than names,
bytes, or hashes alone. It allows arbitrary working-state renaming, any
spanning choice of start root, global left/right reflection, and global
exchange of the tape symbols. The latter two transformations are generous
equivalences for an originality test; global bit exchange also changes the
blank convention. Neither candidate matches any recovered prior
194-state graph under this test. This establishes a structural difference
in the stated corpus, not semantic inequivalence: all correct consistency
machines are intended to have the same blank-halting condition.

\Needspace{9\baselineskip}
\subsection{What the different designs actually change}

\small
\begin{longtable}{>{\raggedright\arraybackslash}p{.18\linewidth}
                  >{\raggedright\arraybackslash}p{.34\linewidth}
                  >{\raggedright\arraybackslash}p{.40\linewidth}}
\toprule Family & Main mechanism & Consequence for this construction\\\midrule\endhead
209--222 and U208 & Unary-power proof encodings, rebracketed decoders,
entry-context and history folds. & These establish the older compiler
lineage. Several 208 reports repeat the same graph; experimental counts
at 207 or below do not establish a smaller calibrated machine.\\
PA196 / PA199 & Diagonal vacuity replaces part of the older logical
interface; separation and axiom heads are reorganized. & The syntactic
calibration, including reverse equality/membership interfaces, is inherited.
Shorter source text need not compile to fewer controls.\\
PA195 & Original bundled, separate-head 195S, cached 195C, combined
diagonal, and repaired-scanner versions. & These vary both command
semantics and decoder allocation. The exact machine identity matters.\\
Primary PA194 & 134 fixed controls plus 67 AX readers minus seven folds.
An alternative allocation uses 68 readers and eight folds. & The two
savings cannot be added; their raw state budgets differ. The present
construction uses the Q family's 67-reader grammar.\\
C / CS / CZ / CSZ & Cached-result MP; cache-only Shift; symmetric
comparison; selected known-zero jumps. & These have different Shift
semantics from QNS. Several C packets duplicate the same table; X and
CSZ also repeat one graph.\\
N / H / HS & No-push MP returns the conclusion in the visible top;
H/HS combine it with a cache-only guard. & QNS retains the original
$t+h=1$ guard. A proof compiler must explicitly reserve the zero that
successful MP no longer pushes.\\
Q / QN & Exchange the two curried premises of the sixth logical head;
QN additionally uses no-push MP. & QN is the exact parent for the main
finite-trace preservation theorem. The head exchange has a fixed
propositional proof, not an assumed logical symmetry.\\
Radix 4R / 5 / 6 and N5/N6 & Change or rearrange the constructor multiplier,
with contextual sharing in its reader; N variants remove the MP push. &
Generic radix recoding must include variable labels. Numerical order is
not preserved by that recoding, so a dominated-axiom translation needs
the supplied unrestricted implicational derivation.\\
J / D / ZR / F / A / R & Known-zero jump substitutions, diagonal-head
parameter changes, register placement, or different AX association. &
Each depends on a concrete source invariant or exact source binding;
none licenses arbitrary branch deletion.\\
SUMS / incoming NS & Reverse the destructive comparison while keeping
the original sum guard; NS also uses no-push MP. & These supply the
symmetric guard transformation used here. Their AX emitter is the
original one, whereas QNS uses the exchanged-premise Q emitter.\\
HA / HAZ & Cached-result MP with the original sum guard, A/D physical
exchange, and optionally a final known-zero AX decrement. & These are
not the no-push visible-top protocol of QNS. The auxiliary MPH itself
already occurred in an earlier N packet.\\
QNS / QS, this report & Q emitter plus the symmetric original-sum guard,
with and without no-push MP. & The compatible composition gives new
literal graphs, removes three source sites from the respective Q
parents, and retains 194 working states.\\
\bottomrule
\end{longtable}
\normalsize

\paragraph{Naming collision.}
The name PA194NS denotes two different supplied designs: an older
cache-only-guard machine that duplicates HS, and a newer incoming
original-sum-guard machine. The comparison records their source paths and
hashes instead of treating the name as an identity.

\paragraph{Formal status.}
At the pinned commit, the top-level PA README explicitly says that the
machine-specific iPRA consistency equivalences for its primary PA194,
PA195, and PA196 targets remain unfinished. It also records substantial
Lean work on literal finite runs, source coverage, clocks, and calculus
composition. This report preserves the distinction: its PA theorem is an
ordinary mathematical proof using the supplied syntactic calibration;
it is not a claim to have closed that separate formalization project.

\section{Exact construction and state accounting}

Start with the native PA194QN table in the incoming QN packet. The new
primary is obtained by the following complete patch; every other row is
unchanged:
\begin{center}
\begin{tabular}{rrll}
\toprule Control & Read & Old action & New action\\\midrule
28&0&$0\,R\,41$&$0\,R\,330$\\
38&1&$1\,R\,20$&$1\,R\,30$\\
41&0&$0\,R\,24$&deleted\\
41&1&$0\,R\,25$&deleted\\
70&0&$0\,R\,41$&$0\,R\,330$\\
310&1&$1\,R\,62$&$1\,R\,330$\\
330&0&absent&$0\,L\,31$\\
330&1&absent&$1\,L\,46$\\
\bottomrule
\end{tabular}
\end{center}
The existing native QN row $(164,1)=(1,R,114)$ remains the no-push
MP branch. Reconstructing the grammar and applying all folds yields the
same table as this direct patch; the package checks both descriptions.

Before the additional controller fold, the new trailer adds one raw
control. The accounting is
\[
 \underbrace{135}_{\text{fixed controls}}
 +\underbrace{67}_{\text{raw AX readers}}
 -\underbrace{7}_{\text{AX folds}}
 -\underbrace{1}_{\text{controller fold}}
 =194.
\]
The raw machine therefore has 202 controls. The final fold merges native
41 into 330. Its validity depends on their different incoming history
bits: each incoming transition to 41 has just written zero and moved
right, while each incoming transition to the raw 330 has written one and
moved right. The shared state can inspect that immediately preceding
tape bit by moving left and then use existing readers 31 and 46.

All eight folds are verified in the final simultaneously folded table.
For each pair there are two possible incoming history bits and two
current read bits, hence four local cases, or 32 in total. Each case
returns to the same tape, head position, and aliased successor as the
original transition after finitely many positive steps. Checking each
fold in isolation would be insufficient because another fold could
change a helper on its simulation path.

The Q emitter itself is unchanged from the supplied QN construction:
190 macro tokens, with 67 distinct raw reader controls. The final
source changes are precisely
\begin{align*}
 \texttt{node168}&:\ (\texttt{node37},\texttt{node57})
       \longmapsto(\texttt{node37},\texttt{node56}),\\
 \texttt{node171}&:\ (\texttt{node170},\texttt{node23})
       \longmapsto
       (\texttt{node170},(\operatorname{inc}H,\operatorname{continueTo}3)).
\end{align*}
The inherited no-push splice is
$\texttt{node163}=(\texttt{node6},\texttt{node54})$.
Both the fully expanded sites and every resolved jump target are supplied;
macro-word equality alone is not used to infer control-flow equality.

\section{Numerical calculus and the inherited PA calibration}
\label{sec:logical}

This section specifies the retained mathematical interface and explains
which arguments are inherited. The detailed calibration texts are
preserved unchanged under \file{references/logical/}, with source hashes
in \file{PROVENANCE.json}.

\subsection{Formula codes and axioms}

Let $p(a,b)=(a+b)(a+b+1)/2+b$ be Cantor pairing and define
\[
 I(A,B)=4p(A,B),\quad M(a,b)=4p(a,b)+1,\quad
 U(a,B)=4p(B,a)+2,\quad E(a,b)=4p(a,b)+3.
\]
Zero denotes truth, and one denotes the target contradiction $M(0,0)$.
The exceptional alias $I(0,0)=0$ is handled explicitly; numerical
formulas are not assumed to be a free term algebra at that point.
Write $\neg B=I(B,1)$ and use the usual classical implicational definitions
of conjunction and existential quantification.

The calculus has truth, the dominated implicational schema
\[
 \rho(a,b,c,u)=I(I(I(c,u),a),I(I(a,c),I(b,c))),\qquad u\le a,
\]
the diagonal schema $\Gamma(B)=I(B,U(B,B))$, the guarded separation
schema below, six ordinary logical heads, and the rules
\[
 \frac{A\qquad I(A,B)}{B}\quad\text{and}\quad
 \frac{B}{U(i,B)}.
\]
All free numeric parameters in the schemas range over natural numbers;
generalization is a rule on theorems.

The six ordinary heads, in the parameters used by the emitter, are
\begin{align*}
 L_1(D,H)&=I(\exists x_D\,U(D,H),H),\\
 L_2(D,A,H)&=I(U(D,U(A,H)),U(A,U(D,H))),\\
 L_3(D,A,H)&=I(M(H,D),I(E(H,A),M(A,D))),\\
 L_4(D,A,H)&=I(E(H,D),I(M(A,H),M(A,D))),\\
 L_5(D,H)&=\exists x_D E(H,D),\\
 L_6^{Q}(D,A,H)&=I(U(D,I(A,H)),I(U(D,A),U(D,H))).
\end{align*}
The last is the exchanged-premise version. The old head is
$I(P,I(Q,R))$, and the new head is $I(Q,I(P,R))$, where
$P=U(D,A)$, $Q=U(D,I(A,H))$, and $R=U(D,H)$.
The retained K/S/MP certificate proves the 61-line exchange template.
K and S are inherited derived rules of the unchanged implicational
basis; neither direction of exchange relies on the quantifier head that
is being replaced. Substituting formulas into that fixed derivation
gives primitive-recursive proof translations in both directions.

For separation put
\begin{align*}
 V(a,b,c)&=I(M(a,b),U(c,M(a,b))),\\
 P(a,b,c)&=I(c,M(c,b)),\\
 B(a,b,c)&=U\bigl(b,
 I(M(a,b),P(a,b,c))\wedge I(M(c,b),c)\bigr),\\
 Q(a,b,c)&=U(c,I(B(a,b,c),\neg P(a,b,c))).
\end{align*}
The guarded head is $G(a,b,c)=\neg I(V(a,b,c),Q(a,b,c))$
when $c\ne a,b$, and is zero otherwise. Formula-code bounds ensure that
the diagonal label $B$ is absent from the formula with code $B$; the
truth-code case is treated separately.

\subsection{Commands and a complete proof-search interface}

A command store is $(t,w,d,h)$: visible theorem, encoded stack, offset,
and cached theorem. Set $[a;w]=p(a,w)$. Its total inverse projections
satisfy $\operatorname{hd}(0)=\operatorname{tl}(0)=0$ and $[0;0]=0$.
This prevents a zero push from being mistaken for a distinct free-list
cell. The primary commands are
\[
 \Shift(t,w,d,h)=
 \begin{cases}\mathsf{halt},&t+h=1,\\
 (h,\operatorname{tl}(w),d,\operatorname{hd}(w)),&t+h\ne1,
 \end{cases}
\]
and, with $b=d+x$,
\[
 \MP_x(t,w,d,h)=
 \begin{cases}(b,w,b,0),&I(t,b)=h,\\
 (0,w,b,0),&I(t,b)\ne h.
 \end{cases}
\]
In particular, success leaves offset $b$, not zero. Neither branch pushes
onto the stack.

For $D=d+x$, $H=h+y$, the AX command returns
\[
 \AX_{A,x,y}(t,w,d,h)=
 \bigl(G(H,D,A),[q_1,\ldots,q_9,t;w],0,0\bigr),
\]
where $(q_1,\ldots,q_7)=(\Gamma(A),L_6^Q,L_5,L_4,L_3,L_2,L_1)$,
$q_8=\rho(D,A,H,d)$ and $q_9=U(d,h)$. The first eight generated
stack entries and the gate are divisible by four. The universal entry
is congruent to two modulo four and is at least two. Thus none of these
generated entries is code one, without assuming consistency.

With $A_0=\AX_{0,0,0}$, the inherited no-push proof compiler uses
\[
 \mathsf{Push}=A_0;\Shift^9;\MP_0,\qquad
 \mathsf{Pop}=A_0;\Shift^{11},\qquad
 \mathsf{Discard}=\MP_0.
\]
For arbitrary natural stores the exact relevant contracts are
\begin{align*}
 \mathsf{Push}(t,w,d,h)&=(0,[t;w],0,0),\\
 \mathsf{Pop}(t,[a;w],d,h)&=(t,w,0,a),\\
 \mathsf{Discard}(t,w,0,h)&=(0,w,0,0).
\end{align*}
Pop is total even when $t=1$, since its last guard tests $U(d,h)+t\ge2$.
Discard works on both success and failure of MP because their outcomes
coincide at zero offset and parameter. The proofs use the literal
Cantor projection identities, including the zero alias.

A proof word $C(p)$ has the entry invariant $d=0$; the initial $A_0$ in
the all-store refutation word below establishes this invariant. A
modus-ponens proof compiler reserves a zero below its two premises:
\[
 \mathsf{Discard};\mathsf{Push};C(I(a,b));\mathsf{Push};
 C(a);\mathsf{Push};\Shift;\Shift;\MP_b;\mathsf{Pop}.
\]
If the recursive premise computations do not halt early, the last Pop
consumes the reserved zero and clears the retained offset, returning
$(b,w,0,0)$. The supplied axiom selectors and generalization words then
give, by induction on a derivation, a word $C(p)$ that either halts or
returns its conclusion with the incoming stack restored and $d=h=0$.
For a refutation,
$A_0;C(p);\Shift$ halts from every incoming store.

This all-store assertion is necessary because the scheduler carries its
store between candidates. It enumerates a cofinal family of fresh
candidates: prefixing the universally halting word by $\AX_{k,0,0}$ gives
unbounded Cantor codes. Consequently starting the fresh counter at two
does not lose completeness. Every decoded continuation decreases and
each source macro terminates. The supplied compiler makes these
finite-word and finite-trace translations primitive recursive.

\subsection{Why this is a PA calibration}

For the soundness direction, interpret membership by finite binary
membership $\operatorname{Bit}(x_a,x_b)$, equality by arithmetic equality,
and $U$ by universal quantification. PA proves that a number is not a
member of itself in this interpretation. The supplied separation proof
uses finite binary comprehension for each substituted arithmetic formula.
It is a primitive-recursive supplier of formal PA proofs, not a uniform
truth predicate or a reflection principle. Coincident labels in the
separation matrix and the alias $I(0,0)=0$ are handled explicitly.

Annotate the visible top, cached head, and canonical stack entries by
their derivations. AX, MP, and Shift preserve those annotations. If the
sum guard halts, one of the two annotated codes is one. This yields a
PA refutation by syntactic translation.

For completeness, the inherited diagonal/separate-head calibration
reconstructs the reverse membership-rewriting and neighborhood-equality
interfaces used by the supplied PA196 proof compiler. This avoids
silently reusing the equality-vacuity axiom removed from that calculus.
The Q exchange template then transports the proof to the exchanged-head
calculus, and the all-store compiler above turns it into a source halt.

The detailed arguments are retained in
\file{separate_head_PA_calibration.tex},
\file{diagonal_calculus.tex}, \file{separate_head_calculus.tex},
\file{pa194q_calibration.tex}, and \file{pa194n_command_proof.tex}
under \file{references/logical/}. The new proof below changes none of
these logical interfaces.

% Standalone article fragment. Needs amsmath, amssymb, amsthm, hyperref.
% The parent article should define theorem and lemma environments.
\section{An independent source-level preservation proof}
\label{sec:source-preservation}

The proposed PA194QNS source is compared here with the exact PA194QN source
from the incoming packet.  This comparison concerns resolved primitive
instructions, not merely equal lists of macro names.  The independent checker
\texttt{check\_source\_preservation.py} imports no builder or physical symbolic
runtime.  It checks that the first 3649 instructions agree, that the old final
five instructions and the new final two instructions are precisely those below,
and that all entries to and exits from the changed region have been accounted
for.  The proof for unbounded registers is the induction given in this section;
the checker's bounded clock diagnostics are not used as a substitute for it.

Logical registers are numbered
$(N,T,W,A,D,H,X,Y,L)=(0,1,2,3,4,5,6,7,8)$.
The complete Shift block begins at source ordinal 3642.
Its only external incoming edge is $48\to3642$.  Both versions share the
following prefix:
\[
 \texttt{while }T>0:\ T:=T-1;\ X:=X+1;
 \qquad
 \texttt{while }H>0:\ H:=H-1;\ T:=T+1;\ X:=X+1.
\]
The old trailer at ordinals 3649--3653 is
\[
 \texttt{while }X>0:\ X:=X-1;\ H:=H+1;
 \qquad X:=X+1;\qquad \texttt{goto }0.
\]
The new trailer at ordinals 3649--3650 is
\[
 H:=H+1;\qquad\texttt{goto }0.
\]
The common comparator occupies source ordinals 0--10, has no incoming edge
from outside except these two trailers, and has only the nonhalting exit
$9\to11$.  Ordinal 10 is the unique halt instruction in either source.
In particular, there is no unchecked jump into the middle of a replacement loop.

\Needspace{12\baselineskip}
\begin{lemma}[Literal destructive comparator]
For arbitrary initial natural values $(H,X)=(u,v)$, the common comparator
reaches source ordinal 10 if and only if $u=v$.  If $u\ne v$, it reaches
ordinal 11.  In either case its final values are $H=X=0$, and every other
register is preserved.  The number of primitive steps to that endpoint is
\[
 C(u,v)=
 \begin{cases}
 3u+2,&u=v,\\
 2u+v+4,&u>v,\\
 3v+2,&u<v.
 \end{cases}
\]
Here arriving at the halt instruction is the endpoint; executing the halt
instruction itself is not counted.
\end{lemma}
\begin{proof}
As long as both counters are positive, the instructions at 0, 1 and 2
decrement each once and return to 0, in three steps.  There are
$\min(u,v)$ such rounds.  At equality the two zero branches reach 10 in
two further steps.  If $u>v$, the next decrement of $H$, the zero test of
$X$, and the jump at 3 lead to the clearing loops in three steps, leaving
$u-v-1$ in $H$.  Clearing it costs $2(u-v-1)$, and the final two zero tests
and exit jump cost three.  The total is $2u+v+4$.
If $u<v$, two steps leave $v-u-1$ in $X$ and reach ordinal 5.  Each unit
remaining in $X$ costs three steps, and the final zero tests and exit cost
three, giving $3v+2$.  These calculations are inductions on natural
counter lengths and exactly match the literal branch targets.
\end{proof}

\begin{lemma}[All-register equivalence of the replacement]
Enter ordinal 3642 with arbitrary natural values $(T,H,X)=(t,h,x)$,
and put $k=t+h+x$.  Both old and new blocks reach ordinal 10 exactly when
$k=1$.  Otherwise they reach ordinal 11 with
\[
                   (T,H,X)=(h,0,0),
\]
and identical values in every other register.
\end{lemma}
\begin{proof}
The first loop has invariant $T+X=t+x$ and terminates with $T=0$.
The second loop has invariants $T+H=h$ and $X+H=x+t+h$ and terminates
with $(T,H,X)=(h,0,k)$.  The old trailer therefore enters the comparator
with $(T,H,X)=(h,k,1)$; the new trailer enters it with $(h,1,k)$.
The comparator lemma proves the same halt decision $k=1$ and the same
complete nonhalting register store.  No zero assumption on the incoming
scratch value $x$ is used.  At a clean command boundary $x=0$, both
implement the original Shift guard $t+h=1$.
\end{proof}

The common prefix takes $3t+4h+2$ primitive steps.  The old and new
trailer-plus-comparator costs are respectively $3k+3+C(k,1)$ and
$2+C(1,k)$.  Their difference is
\[
 \Delta(k)=
 \begin{cases}0,&k=0,\\4,&k=1,\\2k+4,&k\ge2.\end{cases}
\]
This is an exact source-level clock comparison for this complete block.
It is not a claim of a uniform improvement in literal Turing-machine time:
the binary decoder paths and tape excursions have their own costs.

\begin{theorem}[PA-provable source halting equivalence]
Starting from the common clean source entry, PA194QNS halts if and only if
the supplied PA194QN source halts.  The two directions have
primitive-recursive transformations of finite halting traces, and PA
proves their correctness.
\end{theorem}
\begin{proof}
Treat a visit to ordinal 3642, followed by its complete trailer and
comparator, as one block.  All other steps are identical, and the edge
audit shows that this decomposition omits no alternative entrance into
an internal state of the changed block.  The block lemma gives identical
normal exits and matching halting endpoints.  Its loops are total, with
the displayed positive arithmetic clocks.  Replace complete visits one
at a time, and use induction on the finite trace decomposition.
The traces can be generated by primitive recursion using the explicit
counter loops.  The source starts at ordinal 25 and is outside the
changed region; thus the same relation applies from the actual clean
initial source configuration.  Every statement uses only finite codes,
primitive-recursive operations and ordinary arithmetic induction.
\end{proof}

\subsection{What remains inherited}
The new replacement does not alter the Q axiom emitter, the no-push
modus-ponens command, the command scheduler, or any numerical axiom.
The exact command store remains $(t,w,d,h)$, with
\[
 S(t,w,d,h)=
 \begin{cases}\mathsf{halt},&t+h=1,\\
 (h,\operatorname{tl}(w),d,\operatorname{hd}(w)),&t+h\ne1,
 \end{cases}
\]
and, for $b=d+x$,
\[
 \mathsf{MP}_x(t,w,d,h)=
 \begin{cases}(b,w,b,0),&I(t,b)=h,\\(0,w,b,0),&I(t,b)\ne h.
 \end{cases}
\]
In particular, the successful MP offset is $b$, not zero.
The supplied QN proof compiler handles this by reserving a zero below
the two premises and using a final Pop.  The Q emitter changes only
$I(P,I(Q,R))$ to $I(Q,I(P,R))$ in the sixth logical head.
The unchanged dominated implication basis derives unrestricted premise
exchange without using the quantifier-distribution head.  Therefore
fixed proof templates translate the Q calculus to the original calculus
and back.  No use of an arbitrary permutation as a logical axiom is made.

The resulting proof chain has three separate links: the inherited
PA/calculus proof translations; the exact new source equivalence just
proved; and the candidate's positive physical source contracts and
complete blank startup.  Their composition gives the ordinary
mathematical statement
\[
 \mathrm{PA}\vdash
 \operatorname{Halt}_{\rm blank}(\mathrm{PA194QNS})
 \leftrightarrow
 \exists p\,\operatorname{Prf}_{\mathrm{PA}}(p,\ulcorner0=1\urcorner).
\]
Negation gives the desired consistency equivalence.  This written proof
and the executable certificates are not a new end-to-end Lean or Rocq
theorem, and they do not assert the repository's separate iPRA endpoint.

\subsection{The QS companion}
The same block proof also applies to PA194QS.  For its parent source,
restore \texttt{node163=(node6,node55)} in the supplied QN source DAG,
leaving the Q emitter unchanged.  This restores the original successful-MP
dummy push.  The independent \texttt{bind\_q\_sources.py} expands the DAG
and resolves jumps by ordinal subtree intervals.  As a positive control
it first reconstructs the supplied 3654-site QN graph exactly; it then
reconstructs the root's QNS and QS graphs exactly.  The Q-push parent has
3673 instructions.  Its first 3668 agree with QS, whose complete Shift
block starts at 3661 and whose replacement trailer has two instructions.
The only external edge into the block is $48\to3661$; its comparator and
normal exit are the same ordinals 0--10 and 11 as above.  All invariants
and clock formulae are consequently identical after this fixed ordinal shift.

The Q-push source inherits the original separate-head command compiler:
the Q change only replaces the sixth axiom head by its premise-exchanged
counterpart, with the same number and order of stack entries.  The fixed
premise-exchange templates preserve refutations in both directions, while
the original pushing MP and sum-guard Shift are unchanged.  Thus the
same chain of PA-provable finite translations establishes the ordinary
consistency equivalence for QS when composed with its own physical
contracts and startup.  QS and QNS must use their respective original
and no-push proof compilers; their MP commands are not identified with
one another.


\section{From the source to the exact tape machine}
\label{sec:physical}

At a clean source frame the machine is at native control 84. The tape
contains a left marker, the binary source address, a zero gap, a header,
the nine unary register blocks, and a finite nonempty spare list. The
address plus gap length is $48+n$ for arbitrary $n\in\mathbb N$.
The logical-to-physical register map is
\[
 (N,T,W,A,D,H,X,Y,L)\longmapsto(6,2,0,4,3,7,8,5,1).
\]
Every primitive source instruction has a positive-duration physical
contract. Decrement has separate zero and positive contracts. Each
contract preserves the complete tape outside its specified effect and
reaches the exact successor frame, or the literal halt state for a halt
instruction.

Two separate implementations check these obligations. The construction
executor expands the source and uses the inherited register-block
contracts. The literal-tape checker imports neither the compiler nor
those block oracles. It rederives four generic register-list lemmas from
the actual transition rows, resolves the source with two independent
methods, and checks every source outcome over arbitrary natural register
values, both decrement branches, an arbitrary gap parameter, and an
arbitrary finite nonempty list of arbitrary natural spare values.
Its recorded durations have strictly positive constant terms.

The new source-preservation checker imports neither physical executor.
It binds the exact changed instructions, proves there are no extra
entries to the replacement block, and checks the comparator against the
displayed induction. Its finite numerical probes are diagnostics for the
clock expressions; the universal argument is the induction in
Section~\ref{sec:source-preservation}.

\subsection{The complete blank initialization}

From the literal all-zero tape, the primary and companion both reach
source site 25 after exactly 58,004 physical transitions. The resulting
head is at $-2$, control is 84, the logical register vector is
$(2,0,0,0,0,0,0,0,0)$, and the spare list consists of 26 ones. The check
compares every one-cell, every marker, every register block, and the zero
exterior. The parent QN requires 58,154 transitions to the same clean
frame, so this initialization is 150 transitions shorter. This finite
startup comparison is not a claim of a uniform improvement in total
runtime.

\subsection{Both directions of finite-trace transfer}

A finite source halt expands into a finite physical halt by concatenating
the contracts after initialization. Conversely, a deterministic physical
run starting at a clean frame follows the specified contract to its
successor. Each contract is total and has positive duration. Given a
physical first halt at time $T$, iterate the totalized source step
function for $T+1$ instructions. If this bounded iteration finds a source
halt, reflection follows. Otherwise the positive nonhalting simulations,
preceded by startup, produce a deterministic physical trace beyond the
first halt, a contradiction. This is a finite argument inside PA; it
does not require inverse decoding of arbitrary intermediate tapes.
The finite local fold witnesses preserve the argument after all eight
quotients are applied.

The physical symbolic rules, source inductions, and concrete row checks
are ordinary mathematical/executable evidence. They are suitable inputs
for a later formalization; the present release does not label Python
assertions as a proof-assistant kernel theorem. The fresh PA-preservation
argument and the inherited PA-calibration argument are both stated
explicitly, instead of deducing logical equivalence from a long observed
nonhalting run.

\section{Search below 194 and the remaining obstruction}
\label{sec:search}

The achieved state count is 194. No global lower bound is inferred from
the searches. A prior 194-state scanner deletion is explicitly rejected
in the corpus because it misclassifies a legal source address. Another
small diagnostic passes an overly broad source screen but halts during
blank initialization. Such examples motivate checking the exact entry
and complete tape, rather than counting a reachable control graph alone.

One concrete composition attempt combined the cache-only symmetric
guard with the radix-six reader optimization. Its raw allocation was
203 controls. The two proposed extra folds both used native control 41,
so only one could be applied: the result had 195, not 194, controls.
This is a genuine compatibility obstruction for that proposed fold set,
not a lower bound for either semantic protocol.

% Generated by generate_includes.py.
\subsection{Exact QNS word and the combined controller}
The QNS controller was fixed after its symmetric-guard fold, with 134
remaining fixed controls. A 240-second stochastic reassociation run kept
the exact 190-token Q emitter word; parameter permutations were disabled.
It visited 1,366,615 proposal slots, accepted 220,787 moves, and retained
best effective reader cost 60. No zero-penalty hit at cost 59 or below was
found. Such a hit would have been a candidate for 193 states, subject to
full simultaneous-fold and physical verification. Proposal counts are not
counts of distinct grammars, and this is not an exhaustive search.
The source, seed, inputs and actual log are in
\file{research/joint_grammar/}.

\subsection{Larger local windows and source invariants}
Additional experiments use the supplied PA194N6 table, whose input hash
is recorded in \file{research/additional_reductions/RESULTS.json}.
Their endpoints permit arbitrary tape and head effects on halting;
nonhalting endpoints must match exactly. The two operators try a common
replacement action for a pair of controls and a separate replacement
action for each incoming edge of a removed non-self-looping control.
Each local simulation has a 300-step bound.

\begin{center}\small
\begin{tabular}{lrr}
\toprule Quantity & Generic radius 8 & Invariant radius 5\\\midrule
Admissible source outcomes & 4,877 & 4,875\\
Symbolic local views & 100,502 & 21,974\\
Collected tape windows & 642,262 & 29,262\\
Controls covered & 194 & 194\\
Common-action merger candidates & 0 & 0\\
Per-edge bypass candidates & 0 & 0\\
\bottomrule\end{tabular}\end{center}
The generic radius-eight windows cover 17 cells. The source invariant is
an inductive product abstraction of zero, one, and at least two. It
retains 3,775 of 3,776 source sites and removes two impossible outcome
obligations. It proves 3,361 zero and 6,167 positive register coordinates
at the retained sites. An independently checked closure certificate is
included. These additional source facts did not produce a state saving
under either bounded local operator. An earlier halting-aware radius-five
screen over 29,282 windows likewise found no common-action merger.

\subsection{Nested binary tags}
A new proposed encoding replaces $4p$ with $2p$ while retaining distinct
membership, equality and universal constructors; its mathematical
translation is given in Appendix~\ref{sec:nested-code}. The fixed scaffold
shrinks to 133 controls, but the emitted word grows from 190 to 213 tokens.
The 180-second grammar experiment made 633,781 proposals and accepted
124,495. Its best heuristic score was 74: 81 raw readers, minus 11 nominal
folds, plus four scanner penalties. The simultaneous compatibility of
that over-budget grammar was not established. In particular these
numbers do not describe an accepted 203- or 207-state machine. No verified
193- or 194-state candidate resulted from this encoding experiment.

The additional source, contexts, certificates, saved best grammar and
actual logs are included. The finite searches do not exclude longer
replacement traces, stronger relational invariants, different grammars,
different formula encodings, or an unrelated smaller construction.


\section{Transferable ideas and further research questions}

The useful lesson from the ZFC work is to optimize the command interface,
decoder association, and history information together. The no-push MP
protocol and the compiler's reserved-zero technique already have that
lineage, which is credited in the retained N proof. The history-sensitive
fold used here likewise reuses information stored on the tape to avoid
remembering it in two different controls. In particular the supplied N
proof cites the ZFC
\file{candidate-k215/alternatives/zero-spare-transfer} experiment as
ancestry; it is not an uncredited new idea of this report.

ConciseIPRA is relevant to how small proof-checking interfaces and
primitive-recursive proof translations are organized. A smaller iPRA
proof checker does not by itself solve this PA task: substituting it
would require PA-provable translations of refutations in both directions.
No equivalence of $\Con(\mathrm{iPRA})$ and $\Con(\PA)$ is assumed.
The PA calibration used here requires the syntactic supplier of PA
induction proofs for arbitrary-formula finite binary comprehension.
The small-kernel and derived-rule organization in ConciseIPRA is useful
for implementing finite proof translations, but does not replace that
calibration obligation. The detailed dependency review is retained in
\file{references/calibration_audit.txt}.

The following questions remain concrete opportunities:
\begin{enumerate}
 \item Can the symmetric sum controller provide an additional compatible
 AX-reader fold after a different association of the exchanged-head
 emitter? One extra fold, or one fewer raw reader at the same fold count,
 would turn the present allocation into 193 states.
 \item Can a stronger source invariant describing register relationships
 remove a physical branch that is required by the current arbitrary-register
 contracts? Such a reduction needs an inductive reachable-frame invariant,
 not just observations from blank execution. The product zero/one/at-least-two
 invariant in the additional N6 search already eliminates one source
 branch, but did not yield a physical state reduction.
 \item Can a new formula encoding make the constructor multiplier cheaper
 without increasing the AX reader budget? A nested tag encoding with
 $J(a,b)=2p(a,b)$ is a concrete direction, but its full source allocation
 and PA proof translations must both be checked.
 \item Can the two optimizations currently competing for state 41 be
 implemented through a different history context or register placement,
 so that both savings survive in the same table?
 \item Can the all-register comparator theorem be formalized first and
 connected to the existing Lean finite-run bridges, while the logically
 separate calibration-to-PA obligations are discharged explicitly?
\end{enumerate}

\section{Reproduction and artifact map}

Python 3.10 or later and the standard library suffice for construction
and verification. Run with assertions enabled; do not use \texttt{-O} or
\texttt{PYTHONOPTIMIZE}. From the package directory:
\begin{quote}\ttfamily\small
python3 verify\_all.py --quick\\
python3 verify\_all.py
\end{quote}
The quick mode checks reconstruction, source binding, local folds, source
preservation, exchange, exact startup, and corpus novelty. The full mode
additionally reruns both physical symbolic implementations. Fresh replay
outputs go to a temporary directory, so the frozen evidence remains
unchanged. The complete accepted construction is deterministic and does
not require rerunning a stochastic search.

To compile this report, run twice from \file{report/}:
\begin{quote}\ttfamily\small
pdflatex -interaction=nonstopmode -halt-on-error report.tex
\end{quote}

\begin{center}\small
\begin{tabular}{lp{.65\linewidth}}
\toprule Path & Contents\\\midrule
\file{machines/} & Complete dense/native tables, source maps, raw parents,
labels, fold certificates, and construction receipts.\\
\file{verification/} & Independent literal engine, exact source audit,
exchange certificate, and recorded outcomes.\\
\file{comparison/} & Full packet ledger, current census, recovered
frontier tables and exact novelty checks, and the original census.\\
\file{references/} & Inherited mathematical calibration, proof-interface
sources, and the repository's formal-status statement.\\
\file{research/} & Bounded searches, input identities, actual run logs,
and explicit negative-result scope.\\
\file{inputs/}, \file{vendor/} & Frozen construction inputs and attributed
compiler/runtime sources.\\
\file{PROVENANCE.json} & Snapshot, source identities, and attribution.\\
\file{SHA256SUMS} & Hash manifest for every released file except itself.\\
\bottomrule
\end{tabular}
\end{center}

\begin{thebibliography}{9}
\raggedright
\bibitem{repo} Vladimir Reshetnikov and ProveIt contributors,
\emph{ProveIt}, snapshot
\texttt{317b07268681b44ff352d173a162e2a8c241b2fe}.
\url{https://github.com/VladimirReshetnikov/ProveIt}.
\bibitem{primary} \emph{A 194-state Turing machine calibrated to Con(PA)},
\file{Computability/BusyBeaver/Research/PA/docs/pa194_conpa_2026-11-13/}.
\bibitem{qn} \emph{PA194QN research}, incoming archive
\file{docs/incoming/pa194qn_research_2026-11-13.zip} at the pinned snapshot.
\bibitem{sums} \emph{PA194NS and PA194SUMS}, incoming archive
\file{docs/incoming/PA194NS_research.zip} at the pinned snapshot.
\bibitem{q} \emph{PA194Q and PA194CSZ},
\file{Computability/BusyBeaver/Research/PA/docs/pa194q_conpa_2026-11-13/}.
\bibitem{n} \emph{PA194N},
\file{Computability/BusyBeaver/Research/PA/docs/pa194N_conpa_2026-11-13/}.
\bibitem{chang} James Chang, imported PA196 and PA209 developments,
\file{Computability/BusyBeaver/Research/PA/jamesyc-upstream/}.
Original licenses and attribution are preserved in the package.
\end{thebibliography}

\appendix
\clearpage
\section{Complete primary transition table}\label{app:table}
An action $wDq$ means write bit $w$, move $D\in\{L,R\}$, and enter
control $q$. Control 194 is terminal; start is control 0 on the all-zero tape.
The two blocks of columns are independent successive ranges of working controls.
% Generated by generate_includes.py.
{\small\ttfamily
\setlength{\tabcolsep}{14pt}
\begin{longtable}{rrr@{\hspace{2.5em}}rrr}
\toprule q & read 0 & read 1 & q & read 0 & read 1\\\midrule\endfirsthead
\toprule q & read 0 & read 1 & q & read 0 & read 1\\\midrule\endhead
\bottomrule\endfoot
0 & 1R1 & 1R2 & 97 & 0L45 & 0R108\\
1 & 1R3 & 1R1 & 98 & 0L44 & 1R113\\
2 & 0R0 & 1R2 & 99 & 0L114 & 1R115\\
3 & 0R4 & 0R1 & 100 & 0R116 & 0L14\\
4 & 1R5 & 1L6 & 101 & 0R109 & 1R101\\
5 & 1L7 & 1L8 & 102 & 0L45 & 0R101\\
6 & 0L7 & 1L6 & 103 & 0L114 & 1L33\\
7 & 0L9 & 1L6 & 104 & 0R66 & 0R117\\
8 & 0L10 & 0L11 & 105 & 0R56 & 0R117\\
9 & 1R12 & 0L9 & 106 & 0L62 & 1R118\\
10 & 0L9 & 0L10 & 107 & 0R119 & 1R120\\
11 & 1L8 & 1L11 & 108 & 0R14 & 1R108\\
12 & 0R12 & 0R13 & 109 & 0R117 & 1R109\\
13 & 0L14 & 0R13 & 110 & 0R71 & 1R121\\
14 & 1L15 & 0L15 & 111 & 0R122 & 1L123\\
15 & 0L15 & 1L16 & 112 & 0R124 & 1R125\\
16 & 0L17 & 1L16 & 113 & 0R47 & 0L14\\
17 & 0L18 & 1L16 & 114 & 0R90 & 1R98\\
18 & 0L19 & 1L16 & 115 & 0R114 & 1L93\\
19 & 0L20 & 1L16 & 116 & 0R126 & 1R126\\
20 & 0L21 & 1L16 & 117 & 0R127 & 1R117\\
21 & 0L22 & 1L16 & 118 & 0R128 & 1R129\\
22 & 0L23 & 1L16 & 119 & 0R32 & 1R64\\
23 & 0R23 & 1R24 & 120 & 0R31 & 1R96\\
24 & 0R25 & 1R26 & 121 & 0R48 & 1R87\\
25 & 0R27 & 1R28 & 122 & 0R56 & 0R127\\
26 & 0R29 & 1R30 & 123 & 0R65 & 1R33\\
27 & 0R31 & 1R194 & 124 & 0R97 & 1R113\\
28 & 0L32 & 1R33 & 125 & 0R78 & 1R77\\
29 & 0R34 & 1R35 & 126 & 0R47 & 1R47\\
30 & 0L33 & 1L36 & 127 & 0R67 & 1R127\\
31 & 0R37 & 1R38 & 128 & 0R130 & 1R131\\
32 & 0R39 & 1R40 & 129 & 0R132 & 1R133\\
33 & 0R41 & 1R42 & 130 & 0R134 & 1R135\\
34 & 0R43 & 1R44 & 131 & 0R136 & 1R137\\
35 & 0R45 & 1R46 & 132 & 0R138 & 1R86\\
36 & 0R47 & 1R48 & 133 & 0R139 & 0L10\\
37 & 0R49 & 1R41 & 134 & 0R95 & 1R123\\
38 & 0R50 & 1R51 & 135 & 0R140 & 1R126\\
39 & 0R52 & 1R53 & 136 & 0R141 & 1R142\\
40 & 0R54 & 1R55 & 137 & 0R143 & 1R144\\
41 & 0R56 & 0R57 & 138 & 0R145 & 1R40\\
42 & 0R58 & 0L14 & 139 & 0R146 & 1R147\\
43 & 0R59 & 1R60 & 140 & 0L148 & 1L143\\
44 & 0R52 & 1R61 & 141 & 0R28 & 1R149\\
45 & 0L62 & 1R63 & 142 & 0R135 & 1R150\\
46 & 0R64 & 1R65 & 143 & 0R151 & 1R152\\
47 & 0R66 & 0R67 & 144 & 0R153 & 1R154\\
48 & 0L64 & 1L68 & 145 & 0L130 & 1R47\\
49 & 0R69 & 1R51 & 146 & 0R31 & 1R50\\
50 & 0L70 & 0L15 & 147 & 0R155 & 1R146\\
51 & 0R71 & 0L10 & 148 & 0R156 & 1R157\\
52 & 0R72 & 1R73 & 149 & 0R158 & 1R159\\
53 & 0R74 & 1R75 & 150 & 0L159 & 1R160\\
54 & 0R76 & 1R77 & 151 & 0R40 & 1L161\\
55 & 0R78 & 1R79 & 152 & 0R162 & 1R163\\
56 & 0L80 & 0R1 & 153 & 0R164 & 1R164\\
57 & 0R81 & 1R57 & 154 & 0R165 & 1R166\\
58 & 0L72 & 1L65 & 155 & 0R121 & 1R111\\
59 & 0R82 & 0L7 & 156 & 0R75 & 1R165\\
60 & 0R83 & 1R84 & 157 & 0R167 & 1R165\\
61 & 0R41 & 1R85 & 158 & 0R168 & 1R169\\
62 & 0R86 & 1R87 & 159 & 0R170 & 1R157\\
63 & 0R88 & 0L71 & 160 & 0R135 & 1R130\\
64 & 0R30 & 1R87 & 161 & 0R138 & 1R135\\
65 & 0R69 & 1R36 & 162 & 0R171 & 1R172\\
66 & 0L80 & 0R2 & 163 & 0R173 & 1R172\\
67 & 0R89 & 1R67 & 164 & 0L174 & 1R175\\
68 & 0R90 & 1R91 & 165 & 0L162 & 1L176\\
69 & 0R56 & 0R81 & 166 & 0R161 & 1R176\\
70 & 0R65 & 1R92 & 167 & 0R95 & 1R177\\
71 & 0R41 & 0L13 & 168 & 0R134 & 1R178\\
72 & 0R82 & 1R68 & 169 & 0R150 & 1R163\\
73 & 0R48 & 1R93 & 170 & 0R179 & 1R180\\
74 & 0R41 & 1R94 & 171 & 0R167 & 1R140\\
75 & 0R95 & 1R96 & 172 & 0R134 & 1R145\\
76 & 0R97 & 1L98 & 173 & 0R181 & 1R40\\
77 & 0R48 & 1R76 & 174 & 0R175 & 1R182\\
78 & 0R92 & 1R99 & 175 & 0R167 & 1R145\\
79 & 0R48 & 1R100 & 176 & 0R183 & 1R173\\
80 & 1R80 & 1L15 & 177 & 0R86 & 1R61\\
81 & 0R101 & 1R81 & 178 & 0R184 & 1R165\\
82 & 0L45 & 0R89 & 179 & 0R134 & 1R181\\
83 & 0R102 & 1R93 & 180 & 0R185 & 1R185\\
84 & 0R102 & 1R103 & 181 & 0R167 & 1R151\\
85 & 0R104 & 0L14 & 182 & 0R151 & 1R130\\
86 & 0R105 & 1L70 & 183 & 0R164 & 1R186\\
87 & 0R103 & 0L14 & 184 & 0R140 & 1R187\\
88 & 0R106 & 1R107 & 185 & 0R188 & 1L189\\
89 & 0R108 & 1R89 & 186 & 0R190 & 1R40\\
90 & 0R66 & 0R109 & 187 & 0R32 & 1R189\\
91 & 0R110 & 0L87 & 188 & 0R47 & 1R191\\
92 & 0R30 & 1R58 & 189 & 0R148 & 1R192\\
93 & 0R82 & 1L111 & 190 & 0R40 & 1R138\\
94 & 0R99 & 0L14 & 191 & 0R174 & 1R47\\
95 & 0R64 & 1R112 & 192 & 0R182 & 1R193\\
96 & 0R111 & 1R74 & 193 & 0R161 & 1R138\\
\end{longtable}}


\clearpage
\section{The prior 194-state frontier}\label{app:frontier}
The following guide separates accepted manuscript claims, supplementary
designs, and rejected/search objects. Complete hashes, source paths, tables,
and canonical action tuples are in \file{comparison/}. A listed class does
not by itself certify a consistency theorem.
% Generated by generate_includes.py.
{\footnotesize
\begin{longtable}{@{}r >{\raggedright\arraybackslash}p{.21\linewidth}
 >{\raggedright\arraybackslash}p{.14\linewidth}
 >{\raggedright\arraybackslash}p{.54\linewidth}@{}}
\toprule No. & Class guide & Evidence category & Distinguishing mechanism\\\midrule\endhead
1 & Rejected short scanner & Rejected & Seven-zero scanner quotient violates a legal source contract; not an accepted consistency machine.\\[4pt]
2 & PA194-AD & Supported packet & Primary 67/7 decoder with A/D physical register placement; source unchanged.\\[4pt]
3 & Primary PA194 (67/7) & Supported packet & Separate diagonal emitter, original sum-guard controller; 134 fixed + 67 readers - 7 folds.\\[4pt]
4 & PA194 (68/8) & Supported packet & Same separate-diagonal numerical calculus and original sum guard, alternate 68-reader/eight-fold association.\\[4pt]
5 & PA194CS & Supported packet & Cached-result MP and cached-head-only Shift; symmetric destructive comparison removes a transfer loop.\\[4pt]
6 & PA194CZ & Supplementary & Cached controller plus known-zero source jump; a secondary same-size composition.\\[4pt]
7 & PA194C & Supported packet & Cached-result MP and cached-head-only Shift; duplicated as the principal table by five packets.\\[4pt]
8 & PA194-altC & Research seed & Alternate cached decoder in grammar search; not a separate principal verified theorem.\\[4pt]
9 & M6 & Supported packet & Formula-constructor coefficient six, with translated numerical calculus and fresh physical source binding.\\[4pt]
10 & Gamma(D) + CSZA & Supplementary & Gamma(D) emitter with a known-zero cached symmetric controller; only supplementary evidence claimed.\\[4pt]
11 & M4R & Supported packet & Radix-family reassociation with the original coefficient-four resolved source.\\[4pt]
12 & PA194D & Supported packet & Separate diagonal head Gamma(A) changes to Gamma(D); revised selectors and 69-reader/nine-fold realization.\\[4pt]
13 & CSZA final-register 5 & Supplementary & Known-zero final-return sibling; not included among independently replayed main D/J machines.\\[4pt]
14 & PA194J & Supported packet & Ten jumps become known-zero decrements under a finite source invariant; three native rows change.\\[4pt]
15 & M5 & Supported packet & Formula-constructor coefficient five, with explicit logical recoding and physical simulation.\\[4pt]
16 & PA194CSZ = PA194X & Supported packet & Cached symmetric controller with known-zero entry decrement; two names identify the same table.\\[4pt]
17 & CSZA final-register 6 & Supplementary & Known-zero final-return sibling with construction/startup evidence.\\[4pt]
18 & CSZA final-register 3 & Supplementary & Known-zero final-return sibling with construction/startup evidence.\\[4pt]
19 & PA194Q & Supported packet & Exchanges the curried premises in universal distribution, with cached controller; 79 primitive source instructions change.\\[4pt]
20 & Rejected node1-X & Rejected & Known-zero-looking patch is rejected; source correctness alone does not ensure blank initialization.\\[4pt]
21 & C-I-Z & Research seed & Auxiliary cached-source/zero-return search input; no principal release claim.\\[4pt]
22 & PA194F & Supported packet & Cached symmetric AX final jump becomes a decrement of X, proved zero by a four-site cut.\\[4pt]
23 & PA194R & Supported packet & F zero-closing source with reassociated 67-reader grammar; 1111 source addresses change.\\[4pt]
24 & PA194N & Supported packet & Successful MP omits dummy-zero push, returns visible consequent (b,w,b,0), and retains sum guard.\\[4pt]
25 & PA194A & Supported packet & Entry jump becomes a decrement of a source-proved zero register; distinct complete initialization.\\[4pt]
26 & Rejected comparator return 103 & Rejected & Source-equivalent patch passes source contracts but halts after 1361 blank-start steps.\\[4pt]
27 & MPH auxiliary & Research seed & MP returns its result in the cache while retaining the original sum guard; incoming HAZ credits this prior table.\\[4pt]
28 & CS both-zero6 & Research seed & Auxiliary simultaneous zero-return search input; no principal consistency proof claim.\\[4pt]
29 & PA194ZR & Supported packet & Eight comparator returns and AX exit use zero-decrement return; independent 129-site invariant.\\[4pt]
30 & Rejected dec6 & Rejected & Rejected structural return/decrement candidate; do not count as a consistency construction.\\[4pt]
31 & PA194HS = filed PA194NS & Supported packet & No-push visible-top MP combined with cached-head-only symmetric Shift; revised reserved-zero proof compiler.\\[4pt]
32 & PA194N6 & Supported packet & Coefficient-six formula arithmetic plus no-push MP; two packets deliver the same table.\\[4pt]
33 & PA194H & Supported packet & No-push visible-top MP and cached-head-only Shift before symmetric-guard simplification.\\[4pt]
34 & PA194N5 & Supported packet & Coefficient-five formula arithmetic plus no-push MP.\\[4pt]
35 & PA194SUMS & Supported packet & Original sum-guard controller with reversed destructive comparison operands; preserves primary PA194 finite traces.\\[4pt]
36 & incoming PA194NS (sum guard) & Supported packet & SUMS plus no-push MP. Distinct from filed PA194NS=HS, which uses a cached-only guard.\\[4pt]
37 & PA194HA & Supported packet & Prior MPH cached-result semantics plus A/D physical register placement; retains original sum guard.\\[4pt]
38 & PA194HAZ & Supported packet & HA plus known-zero final AX decrement, established by the separate four-site source invariant.\\[4pt]
39 & PA194QN & Supported packet & Q exchanged-premise emitter plus no-push N with original sum guard and revised proof compiler.\\[4pt]
40 & PA194N-AD (JSON auxiliary) & Auxiliary & No-push N with A/D physical register placement, present only inside the NS audit JSON; no primary proof claim.\\[4pt]
\bottomrule\end{longtable}}


\clearpage
\section{All 54 delivery packets}\label{app:ledger}
This ledger records each delivery's stated contribution. Repeated machines
remain visible here as repeated deliveries; the graph census deduplicates them.
The accompanying JSON/CSV gives exact repository paths and wrappers.
% Generated by generate_includes.py.
{\footnotesize
\begin{longtable}{@{}r >{\raggedright\arraybackslash}p{.36\linewidth}
 >{\raggedright\arraybackslash}p{.55\linewidth}@{}}
\toprule No. & Exact delivery name & Stated contribution\\\midrule\endhead
1 & \path{pa_unary_power_cover} & Unary power covering at PA strength (report of 27 October 2026). The U208 packets carry it as a baseline.\\[3pt]
2 & \path{U208_ConPA_Conditional_Context_2026-10-29} & U208 by a conditional tape-context reduction. States a 208-state candidate machine/u208.tm.\\[3pt]
3 & \path{U208_ConPA_Entry_Context_2026-10-29} & U208 by eliminating a state using entry context. States a 208-state machine machine/u208.tm.\\[3pt]
4 & \path{U208_ConPA_History_2026-10-29} & U208 by a finite-history decoder fold. States a 208-state candidate machine/u208.tm.\\[3pt]
5 & \path{U208_ConPA_2026-10-29} & U208 by a three-cell fold. States a 208-state candidate machine/u208.tm.\\[3pt]
6 & \path{u208_research_2026-10-29} & U208 research release. States two 208-state machines, machine/u208.tm and machine/variant\_plain/u208.tm.\\[3pt]
7 & \path{Below_U208_Research_2026-10-29} & Research toward fewer than 208 states. States that no machine with 207 states or fewer was obtained.\\[3pt]
8 & \path{u208_comparison_research_2026-10-29} & Comparison of the 208-state proposals and a search below 208. States that the target was not achieved.\\[3pt]
9 & \path{ConPA_below208_research_2026-10-29} & Comparative audit and reduction search below 208. States that it does not establish a 207-state machine.\\[3pt]
10 & \path{pa_consistency_research_2026-10-29} & Audit and new source reductions. States that no validated machine below 208 states was obtained.\\[3pt]
11 & \path{pa208_research_2026-10-29} & Comparison and research below 208 states. States that it does not establish a valid machine below 208.\\[3pt]
12 & \path{pa_tm_research_2026-11-13} & Comparison and reduction research. States a new 196-state variant with shorter source paths and nothing below 196.\\[3pt]
13 & \path{pa_consistency_research_2026-11-13} & Comparison and diagonal-vacuity research. States a shorter axiom program with a 199-state implementation and nothing below 196.\\[3pt]
14 & \path{conpa_tm_comparison_and_reduction_search} & Comparison and reduction search. States two experimental 196-state command variants and nothing below 196.\\[3pt]
15 & \path{pa196_research_2026-11-13} & PA196 comparison, source factorization and proof obligations. States a source tree of 157 internal nodes and nothing below 196.\\[3pt]
16 & \path{pa_consistency_research} & Research and verification packet. States that no machine below 196 states has been validated.\\[3pt]
17 & \path{pa_tm_comparison_and_reductions} & Comparison and new rewrites, a research snapshot.\\[3pt]
18 & \path{pa_tm_extended_research_2026-11-13} & Comparison and refined reduction search. States that the smallest accepted count is 196.\\[3pt]
19 & \path{conpa_tm_comparison_and_search} & Comparison and a search below 196. States that none was obtained; files named 195 are rejected or diagnostic candidates.\\[3pt]
20 & \path{pa_consistency_reduction} & Comparison and reduction experiments. States two new 196-state variants and a shorter source program with a 203-state implementation.\\[3pt]
21 & \path{pa_consistency_comparison_2026_11_13} & Comparison and new reduction experiments. States that nothing below 196 was obtained.\\[3pt]
22 & \path{PA_consistency_TM_research_2026-11-13} & Comparison and new source reductions. States that no machine below 196 is claimed.\\[3pt]
23 & \path{pa_consistency_research_2026-11-13_2} & Comparison and search below 196. States two new layouts with 196 and 198 states and a command-level reduction.\\[3pt]
24 & \path{conpa_tm_study_2026-11-13} & Comparison and reduction study. States that nothing below 196 was obtained.\\[3pt]
25 & \path{pa_below196_research} & Comparison and reduction experiments. States that nothing below 196 was obtained.\\[3pt]
26 & \path{pa_consistency_comparison_and_new196_2026-11-13} & Comparison and reduction research. States that it contains no accepted machine below 196.\\[3pt]
27 & \path{pa_tm_research_2026-11-13_2} & Comparison and reduction investigation. States same-size work, failed candidates and nothing below 196.\\[3pt]
28 & \path{pa_consistency_cross_archive} & Cross-archive constructions and reduction search. States that it does not certify a 195-state machine.\\[3pt]
29 & \path{pa195_conpa_2026-11-13} & PA195. States a 195-state machine pa195/pa195.tm and an article with a PA proof of the equivalence.\\[3pt]
30 & \path{pa_consistency_reduction_audit_2026-11-14} & Comparison and reduction audit. States nothing below 196 and no new theorem.\\[3pt]
31 & \path{pa_consistency_research_2026-11-14} & Comparison and NEG18. States an 18-instruction encoded negation with a compiled machine of 197 states.\\[3pt]
32 & \path{PA195_ConPA_research_2026-11-14} & A 195-state machine. States the primary machine\_195S.tm and a second table machine\_195C.tm.\\[3pt]
33 & \path{pa195_followup_research} & Comparison and searches below 195. States a graph-distinct 195-state machine M195S8 and nothing below 195.\\[3pt]
34 & \path{pa194_conpa_2026-11-13} & A 194-state machine. States the table construction/tables/pa194.tm and a proof of the equivalence in PA.\\[3pt]
35 & \path{pa194_conpa_2026-11-13_2} & PA194. States a 194-state machine pa194/pa194.tm and an article with a PA proof of the equivalence.\\[3pt]
36 & \path{pa195_diagonal_combined_2026-11-13} & PA195DC. States a different 195-state machine.\\[3pt]
37 & \path{pa194C_conpa_2026-11-13_2} & PA194C. States a materially different 194-state machine, an article with a PA proof of the equivalence, and nothing below 194.\\[3pt]
38 & \path{pa194_new_designs_2026-11-13} & Two new 194-state machines, PA194C and PA194CS. States that nothing below 194 was obtained.\\[3pt]
39 & \path{pa194c_conpa} & PA194C, a cached-head 194-state machine. States that no machine below 194 was established.\\[3pt]
40 & \path{pa194c_conpa_2026-11-13_1} & PA194C. States a different 194-state machine and an article with a PA proof of the equivalence.\\[3pt]
41 & \path{pa194c_conpa_2026-11-13} & PA194C, a cached-head 194-state machine. States an article with a PA proof of the equivalence and nothing below 194.\\[3pt]
42 & \path{pa194_radix_family_2026-11-13} & New 194-state machines calibrated to Con(PA), a radix family.\\[3pt]
43 & \path{pa194j_pa194d_conpa_2026-11-13} & PA194J and PA194D. States two new 194-state machines.\\[3pt]
44 & \path{pa194q_conpa_2026-11-13} & PA194Q and PA194CSZ. States two new 194-state machines and that no machine below 194 was established.\\[3pt]
45 & \path{PA194R_research_2026-11-14} & PA194R. States a different 194-state machine with a new decoder layout and that no machine below 194 was verified.\\[3pt]
46 & \path{pa194N_conpa_2026-11-13} & PA194N. States a new 194-state machine, an article with a PA proof of the equivalence, and nothing below 194.\\[3pt]
47 & \path{pa194zr_conpa_2026-11-13} & PA194ZR. States a new 194-state machine machine/pa194zr.tm, a PA proof of the equivalence, and nothing below 194.\\[3pt]
48 & \path{pa194hs_conpa_2026-11-13} & PA194HS and PA194H. States two 194-state machines and that no machine below 194 was established.\\[3pt]
49 & \path{pa194n5_n6_conpa_2026-11-13} & PA194N5 and PA194N6. States two distinct 194-state machines and claims no machine below 194.\\[3pt]
50 & \path{pa194ns_conpa_2026-11-13} & PA194NS. States a different 194-state machine and that no machine below 194 was established.\\[3pt]
51 & \path{pa194n6_research_2026-11-13} & PA194N6. States a different 194-state machine and that no machine below 194 was verified.\\[3pt]
52 & \path{PA194NS_research.zip} & PA194SUMS reverses the destructive comparison operands while preserving the original sum guard; sum-guard PA194NS adds no-push MP. Both 194 states. A distinct auxiliary PA194N\_AD is supplied only in audit JSON; no sub-194 result.\\[3pt]
53 & \path{pa194haz_research_2026-11-14.zip} & PA194HA and PA194HAZ: cached-result MPH, original sum guard, A/D physical placement; HAZ also changes the final AX jump to a known-zero decrement. Both 194 states. No sub-194 result.\\[3pt]
54 & \path{pa194qn_research_2026-11-13.zip} & PA194QN combines the premise-exchanged Q emitter with no-push N and its original sum guard. 194 states. Two 195-state MP-tail experiments have physical evidence but no completed new consistency proofs. No sub-194 result.\\[3pt]
\bottomrule\end{longtable}}

\clearpage
% Standalone research fragment. Requires amsmath, amssymb, amsthm.
\section{A total code using only twice the Cantor pairing}
\label{sec:nested-code}

The following is a mathematical encoding proposal, not a state-count
improvement established by this packet.  It removes multiplication by four
from formula construction without deleting primitive equality or merging
incompatible formula types.  Its first bounded grammar experiment did not
beat 194 states.

Let $p(a,b)=(a+b)(a+b+1)/2+b$ and $J(a,b)=2p(a,b)$.  Use
\begin{align*}
 I_2(A,B)&=J(A,B),\\
 U_2(a,B)&=J(B,a+1)+1,\\
 M_2(a,b)&=J(J(a,b),0)+1,\\
 E_2(a,b)&=J(J(a,b)+1,0)+1.
\end{align*}
Keep truth code zero, false code one, and unscaled Cantor pairing for the
proof stack and command-word scheduler.  In particular
$M_2(0,0)=1$, $E_2(0,0)=3$, $U_2(0,0)=5$ and $I_2(0,0)=0$.

\begin{lemma}[Complete, unambiguous decoding]
Every nonzero even number is uniquely an implication.  Every odd number
is uniquely a universal formula, membership atom or equality atom.
Formula children and all displayed variable labels are smaller than the
whole nonzero formula code.  The sole truth/constructor alias is
$I_2(0,0)=0$.
\end{lemma}
\begin{proof}
An even number $n>0$ gives its unique pair by unpairing $n/2$.
For an odd $n$, unpair $(n-1)/2$ as $(r,s)$.
If $s>0$, decode $U_2(s-1,r)$.  If $s=0$ and $r$ is even, unpair
$r/2$ as $(a,b)$ and decode $M_2(a,b)$.  If $s=0$ and $r$ is odd,
unpair $(r-1)/2$ and decode $E_2(a,b)$.  These cases are exhaustive
and disjoint.  The inverse equations follow from the two Cantor
projections.  The strict bounds follow from $a,b\le p(a,b)$ and
the outer factor two or outer addition of one.  The cases $n=1,3$
are included and have only label zero.
\end{proof}

\subsection{A primitive-recursive bijection of all formula codes}
Write $I_4,M_4,U_4,E_4$ for the old four-residue constructors.  Define
$F(0)=0$ and, on each nonzero old code, recursively apply the new
constructor to the images of all its arguments.  Variable labels are
recoded by the very same $F$:
\begin{align*}
 F(I_4(A,B))&=I_2(F(A),F(B)),\\
 F(M_4(a,b))&=M_2(F(a),F(b)),\\
 F(U_4(a,B))&=U_2(F(a),F(B)),\\
 F(E_4(a,b))&=E_2(F(a),F(b)).
\end{align*}
The first identity also holds at $A=B=0$ by the truth convention.
The old and new strict-component bounds give course-of-values primitive
recursions for $F$ and its converse $G$, which decodes a new constructor,
recursively applies $G$ to its arguments, and emits the old constructor.
Induction on the natural input code proves
\[
                         G(F(n))=n,\qquad F(G(n))=n.
\]
Thus this is a bijection of all natural codes, rather than an embedding
into a sparse well-formed subset.  It fixes zero and one.  It is not
monotone: $F(2)=5$ and $F(4)=2$.

The simultaneous formula-and-label recoding is essential.  In particular
it makes the diagonal axiom obey the exact identity
\[
 F\bigl(I_4(B,U_4(B,B))\bigr)
       =I_2(F(B),U_2(F(B),F(B))).
\]
Injectivity preserves both label-inequality guards of separation, and
every other logical head commutes with the constructor translation.
Modus ponens and generalization likewise translate literally.

The numerical dominance condition on the restricted $\rho_8$ head does
not survive this recoding.  At such a leaf use the supplied
\texttt{rho8All} template, which derives the unrestricted schema from
MP, the restricted head, reflexivity of the bound, and the inequalities
$a,b\le I(a,b)$.  These inequalities hold for $I_2=2p$ as well as
for $I_4=4p$.  The same construction works for the converse translation.
Therefore both numerical calculi have primitive-recursive,
refutation-preserving proof translations with correctness provable in PA.
The existing ordinary PA calibration is inherited through these
translations; no new equality axiom or set-theoretic principle is assumed.

\subsection{Command and emitter consequences}
The output $q_9=U_2(d,h)$ is odd and at least five.  The other eight
generated heads and the separation gate are even, because their outer
constructor is implication or the result is zero.  Thus no generated
head is one.  Push's final MP compares an even implication with an odd
universal code and necessarily fails.  The original-sum-guard and
no-push command-word arguments retain their exact proofs with these
new constructor meanings.  The Cantor stack and its zero alias are
unchanged.

For an emitter token \texttt{Cons} that now computes $J$ and
\texttt{Push} that saves the old top and puts zero in the top, the
constructor suffix replacements are
\[
\begin{array}{lll}
 \text{old membership:}&\texttt{Cons, inc}&
       \texttt{Cons, Push, Cons, inc},\\
 \text{old universal:}&\texttt{Cons, inc, inc}&
       \texttt{inc, Cons, inc},\\
 \text{old equality:}&\texttt{Cons, inc, inc, inc}&
       \texttt{Cons, inc, Push, Cons, inc}.
\end{array}
\]
For the universal replacement the current top is the binder label and
the preceding stack entry is the body, exactly as in the old transposed
constructor.  Its first increment therefore produces $a+1$.
The two atom replacements create a zero right component for the outer
pair while preserving the older stack.

The independent reduction workstream found 31 implication, 15 universal,
10 membership and 3 equality constructor sites in the 190-token emitter.
The replacements give 213 tokens.  Removing the four-increment decoder
gives a theoretical fixed-control budget of 133, so reaching 193 would
require a compatible effective reader cost of at most 60 under that
accounting.

The recorded 180.027-second grammar search made 633,781 proposals and
accepted 124,495 moves.  Its best heuristic score was 74: 81 raw readers,
minus 11 nominal folds, plus four scanner penalties.  The penalties are
part of the search objective, not a count of machine controls, and the
nominal folds were not certified as a compatible simultaneous quotient
at this endpoint.  Consequently neither 203 nor 207 is an accepted
machine count from this run.  The scoring experiment failed to yield
a verified 193- or 194-state candidate.  It establishes no lower bound
on the encoding.  A different emitter organization or jointly optimized
dictionary remains a research possibility.

\end{document}
